\section{Introduction}

Forensic analytics can prioritize investigation. They cannot prove fraud,
certify correctness, or replace risk-limiting audits. This boundary is not a
footnote for the W-series; it is the first claim every method must make. A
score from this paper is a reason to ask a question, not an answer to that
question.

This method paper covers two time- and sequence-oriented signals. The first is
an abrupt shift in a reported-results series: turnout, vote share, margin,
residuals, or other normalized features ordered by reporting time or by batch
sequence. Such a shift may be a change point worth review. The second is a set
of audit discrepancies that cluster by source, scanner, batch, or similar
operational unit instead of scattering as the null model would expect.

Both signals are useful because RCOUNT packages preserve the record structure
needed to replay them: status events, batch manifests, source hashes, and audit
observations. Both are also easy to overread. Mail batches reported first,
late-arriving provisional ballots, one high-volume scanner, or a local reporting
workflow can create legitimate discontinuities and clusters. The purpose here is
therefore narrow: instantiate the W.01 forensic report contract for these two
families and keep the result investigative, reproducible, and separate from
certification evidence.
